\documentclass[12pt]{article}
\usepackage{pgfplots}
\pgfplotsset{compat=1.18}
\usepackage{chemfig}
\usepackage[version=4]{mhchem}
\usepackage[utf8]{inputenc}
\usepackage[french]{babel}
\usepackage{amsmath}
\usepackage{circuitikz}

\begin{document}

\section*{Circuit RC et Charge du Condensateur}

\subsection*{Schéma du Circuit RC}
\begin{center}
\begin{circuitikz}
  % Alimentation
  \draw (0,0) to[voltage source, l=$\mathcal{V}_0$] (0,3) 
        % Résistance
        to[R=$R$, i=$i$] (3,3)
        % Condensateur
        to[C=$C$] (3,0)
        % Fermeture du circuit
        to[short] (0,0);
\end{circuitikz}
\end{center}

\subsection*{Graphique de Charge}
\begin{center}
\begin{tikzpicture}[scale=1.0]
  % Axes
  \draw[->] (0,0) -- (5,0) node[right] {Temps (s)};
  \draw[->] (0,0) -- (0,4) node[above] {Tension (V)};
  
  % Courbe de charge
  \draw[blue, thick, domain=0:5, samples=100] plot ({\x}, {4*(1 - exp(-\x/0.5))});
  
  % Légende
  \node[blue] at (3,3) {$V_C(t) = V_0(1 - e^{-t/RC})$};
  
  % Points importants
  \fill (0.5, 2.3) circle (2pt) node[above right] {$t = RC$};
\end{tikzpicture}
\end{center}

\subsection*{Équation Théorique}
La tension aux bornes du condensateur suit la loi suivante :
$$
V_C(t) = V_0 \left(1 - e^{-\frac{t}{RC}}\right)
$$

\end{document}